\section*{Bab \ref{ch:g5:stars-night-sky} --- Bintang dan Langit Malam}

\begin{solution}{exo:g5:stars-night-sky:1}
\omterm{def:g4:solar-system:star}{Bintang} adalah sebuah matahari: bola raksasa dari zat yang berpijar
panas dan membuat \omterm{def:g1:light-and-shadows:source}{cahayanya} sendiri. Mereka tampak sebagai percikan
hanya karena jaraknya --- \omterm{def:g1:light-and-shadows:source}{cahaya} pun menghabiskan bertahun-tahun di
jalan dari yang terdekat.
\end{solution}

\begin{solution}{exo:g5:stars-night-sky:2}
Ya --- mereka bersinar pada tengah hari persis seperti pada tengah
malam. Langit siang dibanjiri \omterm{def:g1:light-and-shadows:source}{cahaya} matahari yang dihamburkan \omterm{def:g2:air-around-us:air}{udara},
dan tidak ada percikan yang dapat tampil di sebelah silau itu, seperti
tidak ada lilin yang tampil di sebelah api unggun.
\end{solution}

\begin{solution}{exo:g5:stars-night-sky:3}
Yang putih kebiruan --- warna \omterm{def:g4:solar-system:star}{bintang} berjalan dari jingga (paling
sejuk) lewat kuning sampai putih kebiruan (paling panas). Bara yang
berpijar jingga dibandingkan dengan baja yang dipanaskan sampai
membara putih menceritakan kisah yang sama.
\end{solution}

\begin{solution}{exo:g5:stars-night-sky:4}
Pola \omterm{def:g4:solar-system:star}{bintang} terang yang diberi nama, diciptakan sebagai peta langit
sekaligus alat bantu ingatan. \omterm{def:g4:solar-system:star}{Bintangnya} bukan tetangga yang
sebenarnya --- mereka hanya kebetulan berbaris jika dilihat dari Bumi.
\end{solution}

\begin{solution}{exo:g5:stars-night-sky:5}
Temukan rasi Biduk; ambillah kedua \omterm{def:g4:solar-system:star}{bintang} penunjuk di tepi depan
pancinya; panjangkan garisnya kira-kira lima kali jaraknya menyeberangi
langit; \omterm{def:g4:solar-system:star}{bintang} yang dicapai itulah \omterm{met:g5:stars-night-sky:northstar}{Bintang Utara}, dan di bawahnya
pada cakrawala terletak utara.
\end{solution}

\begin{solution}{exo:g5:stars-night-sky:6}
Ia berdiri hampir tepat di atas sumbu putar Bumi, sehingga putaran
langit meninggalkannya nyaris tidak bergerak --- sebuah tanda utara
yang tetap, yang menggantikan \omterm{def:g4:magnets-and-compass:compass}{kompas}.
\end{solution}

\begin{solution}{exo:g5:stars-night-sky:7}
Bumilah yang berputar, membawa kita berkeliling sekali sehari. Putaran
yang sama mengarak Matahari menyeberangi langit siang --- satu
putaran, dua arak-arakan.
\end{solution}

\begin{solution}{exo:g5:stars-night-sky:8}
\omterm{def:g4:solar-system:star}{Bintang} samar yang tak terhitung banyaknya, terlalu berjejal untuk
dipisahkan mata: itulah \omterm{ex:g5:stars-night-sky:milkyway}{Bima Sakti} --- kota \omterm{def:g4:solar-system:star}{bintang} kita sendiri yang
pipih, dilihat dari sisinya, dari dalam.
\end{solution}

\begin{solution}{exo:g5:stars-night-sky:9}
Bumi yang berputar memutar seluruh langit mengelilingi sumbu yang
melewati \omterm{met:g5:stars-night-sky:northstar}{Bintang Utara}. Di dekat titik diam itu, \omterm{def:g4:solar-system:star}{bintang} menunggangi
lingkaran kecil --- kameranya merekam cincin. Ke arah timur,
putarannya membawa \omterm{def:g4:solar-system:star}{bintang} naik melewati cakrawala menyusuri jalan
yang menyerong --- kameranya merekam goresan yang menaik. Satu
putaran, dua cara memandangnya.
\end{solution}

\begin{solution}{exo:g5:stars-night-sky:10}
Arah: penunjuk Biduk memberinya \omterm{met:g5:stars-night-sky:northstar}{Bintang Utara}, dan bersamanya utara
--- tanpa memerlukan jarum. \omterm{ex:g2:measuring-time:clock}{Jam}: Biduk sendiri berputar mengelilingi
\omterm{met:g5:stars-night-sky:northstar}{Bintang Utara} seperti jarum sebuah \omterm{ex:g2:measuring-time:clock}{jam} raksasa, satu putaran penuh
tiap hari; mengamati seberapa jauh ia sudah berayun memberitahunya
bagaimana malam berlalu.
\end{solution}

\begin{solution}{exo:g5:stars-night-sky:11}
Tidak: kamu melihat \omterm{def:g4:solar-system:star}{bintang} itu sebagaimana keadaannya \emph{dahulu},
bertahun-tahun lalu, ketika \omterm{def:g1:light-and-shadows:source}{cahaya} malam ini meninggalkannya --- sang
kurir, bukan sang pengirim. Karena itu setiap teleskop memandang ke
masa lalu: makin jauh \omterm{def:g4:solar-system:star}{bintangnya}, makin tua kabarnya --- dan justru
itulah yang membuat \omterm{def:g1:light-and-shadows:source}{cahaya} yang jauh begitu berharga bagi para
astronom.
\end{solution}
